
% transition.tex
\subsection{11.1 transition 过渡渐变效果}
\begin{itemize}[leftmargin=*]
\item 过渡动画概念：样式状态平滑渐变切换
\item transition-property 指定参与过渡的样式属性
\item transition-duration 设置过渡持续时长
\item transition-timing-function 运动速度曲线
\item transition-delay 过渡触发延迟等待时间
\item 简写复合属性统一配置过渡参数
\item 鼠标悬浮触发样式过渡变化
\item 点击获取焦点触发过渡效果
\item 类名切换驱动动态过渡动画
\item 无过渡样式瞬间切换对比区别
\end{itemize}

\subsection{11.2 运动曲线函数分类}
\begin{itemize}[leftmargin=*]
\item linear 匀速直线运动曲线
\item ease 先慢后快再放缓默认曲线
\item ease-in 由慢到快加速运动
\item ease-out 由快到慢减速运动
\item ease-in-out 两端慢中间快对称曲线
\item cubic-bezier 自定义贝塞尔曲线
\end{itemize}

\subsection{11.3 过渡效果完整代码示例}
\subsubsection{示例1 分离属性写法}
\begin{verbatim}
<!DOCTYPE html>
<html lang="zh-CN">
<head>
    <meta charset="UTF-8">
    <title>分离属性过渡写法</title>
    <style>
        .box {
            width: 100px;
            height: 100px;
            background-color: #409eff;
            /* 单独拆分过渡属性 */
            transition-property: width, background-color;
            transition-duration: 0.5s;
            transition-timing-function: ease;
            transition-delay: 0s;
        }
        .box:hover {
            width: 200px;
            background-color: #f56c6c;
        }
    </style>
</head>
<body>
    <div class="box"></div>
</body>
</html>
\end{verbatim}

\subsubsection{示例2 复合简写写法}
\begin{verbatim}
<!DOCTYPE html>
<html lang="zh-CN">
<head>
    <meta charset="UTF-8">
    <title>过渡简写语法</title>
    <style>
        .box2 {
            width: 100px;
            height: 100px;
            background: #67c23a;
            border-radius: 4px;
            /* 简写格式：属性 时长 曲线 延迟 */
            transition: all 0.4s ease-in-out 0.1s;
        }
        .box2:hover {
            width: 180px;
            height: 180px;
            border-radius: 50\%;
            background: #e6a23c;
        }
    </style>
</head>
<body>
    <div class="box2"></div>
</body>
</html>
\end{verbatim}

\subsubsection{示例3 多种运动曲线对比}
\begin{verbatim}
<!DOCTYPE html>
<html lang="zh-CN">
<head>
    <meta charset="UTF-8">
    <title>运动曲线对比</title>
    <style>
        .wrap {margin: 20px;}
        .curve-item {
            width: 80px;
            height: 60px;
            background: #909399;
            margin: 10px 0;
            transition: margin-left 1s;
        }
        .linear {transition-timing-function: linear;}
        .ease {transition-timing-function: ease;}
        .ease-in {transition-timing-function: ease-in;}
        .ease-out {transition-timing-function: ease-out;}
        .ease-in-out {transition-timing-function: ease-in-out;}

        .wrap:hover .curve-item {
            margin-left: 200px;
        }
    </style>
</head>
<body>
    <div class="wrap">
        <div class="curve-item linear">linear</div>
        <div class="curve-item ease">ease</div>
        <div class="curve-item ease-in">ease-in</div>
        <div class="curve-item ease-out">ease-out</div>
        <div class="curve-item ease-in-out">ease-in-out</div>
    </div>
</body>
</html>
\end{verbatim}

\subsection{11.4 过渡核心属性说明}
\begin{itemize}[leftmargin=*]
\item transition-property：设定参与渐变的样式，all代表全部属性
\item transition-duration：动画执行时长，单位s、ms
\item transition-timing-function：控制动画快慢节奏曲线
\item transition-delay：触发后延迟执行的等待时间
\end{itemize}

\subsection{11.5 过渡使用特点}
\begin{itemize}[leftmargin=*]
\item 常用触发方式：\:hover、\:active、类名切换、JS修改样式
\item 仅支持两个状态之间平滑切换动画
\item 过渡效果不脱离标准文档流，仅视觉形态变化
\end{itemize}
